\chapter{1995年全国卷（文）}
\section{单选题}

\begin{problem}
已知集合 \(I = \{0, -1, -2, -3, -4\}\)，集合 \(M = \{0, -1, -2\}, N = \{0, -3, -4\}\)，则 \(\overline{M} \cap N =\)（\quad）

\choices
    {\(\{0\}\)}
    {\(\{-3,-4\} \)}
    {\(\{-1,-2\} \)}
    {\(\varnothing\)}

\end{problem}

\begin{problem}
函数 \(y = \frac{1}{x + 1} \) 的图象是（\quad）

\bitmapfigure[max width=\linewidth,max totalheight=5.4cm]{img/1995/paper_1995_0178/q2_fig1.png}

\end{problem}

\begin{problem}
函数 \(y = 4\sin \left(3x + \frac{\pi}{4}\right) + 3\cos \left(3x + \frac{\pi}{4}\right) \) 的最小正周期是（\quad）

\choices
    {\(6\pi\)}
    {\(2\pi\)}
    {\(\frac{2\pi}{3} \)}
    {\(\frac{\pi}{3}\)}

\end{problem}

\begin{problem}
正方体的全面积是 \(a^2\)，它的顶点都在球面上，这个球的表面积是（\quad）

\choices
    {\(\frac{\pi a^2}{3}\)}
    {\(\frac{\pi a^2}{2}\)}
    {\(2\pi a^{2}\)}
    {\(3\pi a^{2}\)}

\end{problem}

\begin{problem}
如图，若图中的直线 \(l_{1}, l_{2}, l_{3} \) 的斜率分别为 \(k_{1}, k_{2}, k_{3} \) ，则（\quad）

\bitmapfigure[max width=\linewidth,max totalheight=5.4cm]{img/1995/paper_1995_0178/q5_fig1.png}

\end{problem}

\begin{problem}
双曲线 \(3x^{2} - y^{2} = 3\) 的渐近线方程是（\quad）

\choices
    {\(y = \pm 3x\)}
    {\(y = \pm \frac{1}{3} x\)}
    {\(y = \pm \sqrt{3} x\)}
    {\(y = \pm \frac{\sqrt{3}}{3} x\)}

\end{problem}

\begin{problem}
使 \(\sin x \leqslant \cos x\) 成立的 \(x\) 的一个变化区间是（\quad）

\choices
    {\(\left[-\frac{3\pi}{4}, \frac{\pi}{4}\right]\)}
    {\(\left[-\frac{\pi}{2}, \frac{\pi}{2}\right]\)}
    {\(\left[\frac{\pi}{4}, \frac{3\pi}{4}\right]\)}
    {\([0, \pi]\)}

\end{problem}

\begin{problem}
\(x^{2} + y^{2} - 2x = 0\) 和 \(x^{2} + y^{2} + 4y = 0\) 的位置关系是（\quad）

\choices
    {相离}
    {外切}
    {相交}
    {内切}

\end{problem}

\begin{problem}
已知 \(\theta\) 是第三象限角，且 \(\sin^4\theta + \cos^4\theta = \frac{5}{9}\)，那么 \(\sin 2\theta\) 等于（\quad）

\choices
    {\(\frac{2\sqrt{2}}{3}\)}
    {\(-\frac{2\sqrt{2}}{3}\)}
    {\(\frac{2}{3} \)}
    {\(-\frac{2}{3}\)}

\end{problem}

\begin{problem}
如图，\(ABCD - A_{1}B_{1}C_{1}D_{1}\) 是正方体，\(B_{1}E_{1} = D_{1}F_{1} = \frac{A_{1}B_{1}}{4}\)，则 \(BE_{1}\) 与 \(DF_{1}\) 所成的角的余弦值是（\quad）

\bitmapfigure[max width=0.42\linewidth,max totalheight=5.6cm]{img/1995/paper_1995_0178/q10_fig1.png}

\choices
    {\(\frac{15}{17}\)}
    {\(\frac{1}{2}\)}
    {\(\frac{8}{17}\)}
    {\(\frac{\sqrt{3}}{2}\)}

\end{problem}

\begin{problem}
已知 \(y = \log_a(2 - x) \) 是 \(x \) 的增函数，则 \(a \) 的取值范围是（\quad）

\choices
    {\((0,2)\)}
    {\((0,1)\)}
    {\((1,2)\)}
    {\((2, + \infty)\)}

\end{problem}

\begin{problem}
在 \((1 - x^{3})(1 + x)^{10}\) 的展开式中，\(x^{5}\) 的系数是（\quad）

\choices
    {\(-297\)}
    {\(-252\)}
    {297}
    {207}

\end{problem}

\begin{problem}
已知直线 \(l \perp\) 平面 \(\alpha\)，直线 \(m \subset\) 平面 \(\beta\)，有下面四个命题：

\begin{circlelist}
\item \(\alpha \parallel \beta \Rightarrow l \perp m\)；
\item \(\alpha \perp \beta \Rightarrow l \parallel m\)；
\item \(l \parallel m \Rightarrow \alpha \perp \beta\)；
\item \(l \perp m \Rightarrow \alpha \parallel \beta\).
\end{circlelist}

其中正确的两个命题是（\quad）

\choices
    {\circled{1}与\circled{2}}
    {\circled{3}与\circled{4}}
    {\circled{2}与\circled{4}}
    {\circled{1}与\circled{3}}

\end{problem}

\begin{problem}
等差数列 \(\{a_{n}\}, \{b_{n}\}\) 的前 \(n\) 项和分别为 \(S_{n}\) 与 \(T_{n}\)，若 \(\frac{S_{n}}{T_{n}} = \frac{2n}{3n + 1}\)，则 \(\displaystyle\lim_{n\to \infty}\frac{a_n}{b_n}\) 等于（\quad）

\choices
    {1}
    {\(\frac{\sqrt{6}}{3} \)}
    {\(\frac{2}{3} \)}
    {\(\frac{4}{9} \)}

\end{problem}

\begin{problem}
用 1, 2, 3, 4, 5 这五个数字组成没有重复数字的三位数，其中偶数共有（\quad）

\choices
    {24 个}
    {30 个}
    {40 个}
    {50 个}

\end{problem}

\section{填空题}

\begin{problem}
方程 \(\log_2(x + 1)^2 + \log_4(x + 1) = 5\) 的解是\fillinblank{}.

\end{problem}

\begin{problem}
已知圆台上，下底面圆周都在球面上，且下底面过球心，母线与底面所成的角为 \(\frac{\pi}{3}\)，则圆台的体积与球体积之比为\fillinblank{}.

\end{problem}

\begin{problem}
函数 \(y = \cos x + \cos \left(x + \frac{\pi}{3} \right) \) 的最大值是\fillinblank{}.

\end{problem}

\begin{problem}
若直线 \(l\) 过抛物线 \(y^{2}=4(x+1) \) 的焦点，并且与 \(x\) 轴垂直，则 \(l\) 被抛物线截得的线段长为\fillinblank{}.

\end{problem}

\begin{problem}
四个不同的小球放入编号为 1, 2, 3, 4 的四个盒中，则恰有一个空盒的放法共有\fillinblank{}种（用数字作答）.

\end{problem}

\section{解答题}

\begin{problem}
解方程 \(3^{x + 2} - 3^{2 - x} = 80\).

\end{problem}

\begin{problem}
设复数 \(z = \cos \theta + \i\sin \theta, \theta \in (\pi, 2\pi)\)，求复数 \(z^2 + z\) 的模和辐角.

\end{problem}

\begin{problem}
设 \(\{a_{n}\} \) 是由正数组成的等比数列， \(S_{n} \) 是其前 \(n\) 项和，证明： \(\frac{\log_{0.5}S_{n}+\log_{0.5}S_{n+2}}{2}>\log_{0.5}S_{n+1}. \)

\end{problem}

\begin{problem}
如图，圆柱的轴截面 \(ABCD\) 是正方形，点 \(E\) 在底面的圆周上，\(AF \perp DE \) ，\(F\) 是垂足.

\begin{enumerate}
\item 求证：\(AF \perp DB\)；
\item 如果圆柱与三棱锥 \(D-ABE\) 的体积的比等于 \(3\pi \) ，求点 \(E\) 到截面 \(ABCD\) 的距离.
\end{enumerate}

\bitmapfigure[max width=0.34\linewidth,max totalheight=5.2cm]{img/1995/paper_1995_0178/q24_fig1.png}

\end{problem}

\begin{problem}
某地为促进淡水鱼养殖业的发展，将价格控制在适当范围内，决定对淡水鱼养殖提供政府补贴. 设淡水鱼的市场价格为 \(x\) 元/千克，政府补贴为 \(t\) 元/千克. 根据市场调查，当 \(8 \leqslant x \leqslant 14\) 时，淡水鱼的市场日供应量 \(P\) 千克与市场日需求量 \(Q\) 千克近似满足关系：\(P = 1000(x + t - 8)\) \((x \geqslant 8, t \geqslant 0)\)，\(Q = 500\sqrt{40 - (x - 8)^2}\) \((8 \leqslant x \leqslant 14)\). 当 \(P = Q\) 时市场价格称为市场平衡价格.

\begin{enumerate}
\item 将市场平衡价格表示为政府补贴的函数，并求出函数的定义域；
\item 为使市场平衡价格不高于每千克 10 元，政府补贴至少为每千克多少元？
\end{enumerate}

\end{problem}

\begin{problem}
已知椭圆 \(\frac{x^2}{24} + \frac{y^2}{16} = 1\)，直线 \(l: x = 12\). \(P\) 是 \(l\) 上一点，射线 \(OP\) 交椭圆于点 \(R\)，又点 \(Q\) 在 \(OP\) 上且满足 \(|OQ| \cdot |OP| = |OR|^2\)，当点 \(P\) 在 \(l\) 上移动时，求点 \(Q\) 的轨迹方程，并说明轨迹是什么曲线.
\end{problem}
